\centerline{\bf \Huge Инварианты}
\title{Инварианты}

\begin{flushright}
    \scriptsize
    — Сколько у нас денег?\\
    — Три доллара.\\
    — Блин, Картман, ты ж говорил, у нас куча денег!\\
    — Да. Но я не учел тот факт, что слаб в математике.\\
    Эрик Картман. South Park
\end{flushright}

Инвариант $-$ это некоторая характеристика объекта, которая остается неизменной в результате какого-либо процесса. Инвариант (от латинского invarians, в родительном падеже invariantis) – «неизменяющийся». Термин был введён английским математиком Джеймсом Джозефом Сильвестром в 1851 году.

С помощью инварианта можно показать невозможность или возможность достижения некоторого состояния объекта, а значит, инвариант является основой для решения большого класса задач.
В качестве инвариантов рассматриваются: \\
-чётность (нечётность)\\
-остаток от деления\\
-перестановки\\
-раскраски\\
-сумма или произведение каких-нибудь чисел\\

\problem{1}НеНачальник решил наградить Эльзяту и Германа за шикарное решение задач на логику и захотел подарить каждому по 1000001 юаней! У НеНачальника в обращении есть огромное количество купюр по 312 юаней и 160802 юаней. Сможет ли он подарить Эльзяте и Герману ровно по 1000001 китайских юаней или это все скам?\\

\problem{2}Летом 2023 года Федор смог оседлать гигансткую лягушку и начал прыгать вдоль прямой. Сначала они прыгнули на 1 м, затем на 3 м в том же или в противоположном направлении, затем на 5 м в том же или в другом направлении и т. д. НеНачальник пообещал подарить 1000001 японских йен, если случится так, что Федор с лягушкой окажутся в исходной точке после 993-го своего прыжка. Это снова скам?\newpage

\problem{3}НеНачальник выписал на доску следующие числа: 1, 0, 1, 0, 0, 1, 0, 0. Затем Амуля решила к любым двум записанным числам прибавить одно и то же натуральное число. НеНачальник пообещал, что если у Амули после некоторого количества таких операций получится достичь того, чтобы все записанные числа оказались равными, то он подарит ей 1000001 корейских вон. Подарит ли?\\

\problem{4}Герман получил 4 балла за домашнее задание по комбинаторике и в порыве гнева разорвал листок со своими решениями на десять кусков. Затем один из получившихся кусков он разорвал еще на 10 кусков. НеНачальник пообещал подарить Герману 30 рублей (на хлеб), если по завершении релаксации у него окажется 2026 кусочков, а если останется 2023 кусочка, то 1000001 монгольских тугриков. Сколько в итоге НеНачальник подарит Герману?\\

\problem{5}НеНачальник записал следующие числа в вершинах куба: 2, 0, 0, 3, 1, 9, 5, 7. За один ход разрешается Эвене прибавить к числам, стоящим на концах одного ребра, одно и то же целое число. НеНачальник пообещал подарить Эвене 1000001 бразильских реалов, если во всех вершинах получатся нули. НеНачальник просто рофлит? \\

\problem{6}Егор нарисовал на доске квадрат 5x5 и заполнил его числами так, что произведение чисел в каждой строке отрицательно. НеНачальник пообещал подарить Егору 1000001 казахстанских тенге, если после его заполнения произведение чисел в каждом столбике будет положительным, а если найдется столбик, в котором произведение будет отрицательным, то подарит 17 рублей (на проезд). Сколько в итоге НеНачальник подарит Егору?\\

\problem{7}НеНачальник написал на доске в некотором порядке 2016 плюсов и 2017 минусов. Время от времени Саша подходит к доске, стирает любые два знака и пишет вместо них один, причем если она стерла одинаковые знаки, то вместо них она пишет плюс, а если разные, то минус. После нескольких таких действий на доске остался только один знак. НеНачальник пообещал подарить Саше 50 рублей (на сникерс), если на доске останется минус, а если плюс, то 1000001 турецких лир. Сколько в итоге НеНачальник подарит Саше?\\

\problem{8} НеНачальник создал два печатных автомата. Первый по карточке с числами $(a; \ b; \ c)$ выдает карточку с числами
$\biggl(\dfrac{a+b}{2}; \ \dfrac{b+c}{2}; \ \dfrac{a+c}{2}\biggl)$, а второй по карточке с числами $(a; b; c) -$  карточку с числами $(2a-b; \ 2b-c; \ 2c-a)$. НеНачальник пообещал подарить Бате и Глебу 2000002 мексиканских песо, если, используя эти автоматы, они из карточки (2,8; -1,7; 16) получат карточку (1,73; 2; 0,4). НеНачальник ужасный дядя?!\\


\problem{9} НеНачальник записал на доске три числа: 2011, 2012, 2013. За один ход разрешается заменить числа $a, \ b, \ c$ на $\dfrac{ab}{c},\dfrac{ac}{b},\dfrac{bc}{a}$. НеНачальник пообещал подарить Очиру Владимировичу и Расулу Владимировичу по 1000001 кг чистого золота, если они докажут почему невозможно за любое количество ходов получить числа 2008, 2012, 2016. Таким образом НеНачальник потерял большую часть запасов золота. А получится ли у вас выиграть у НеНачальника золото?